% PREAMBULO DE LA VERSION FINAL 4 (copia del de capitulos_final3/, 2026-10-20).
% Cada version final lleva su preambulo, su bibliografia y sus figuras rehechas dentro de
% capitulos_finalN/, para que corregirla no altere la version anterior.
%
% =====================================================================
%  preambulo.tex  -- Paquetes y configuracion de la tesis EduFEM (final4)
%  Clase: report (12pt, a4paper). Motor: XeLaTeX + biber.
%     latexmk -xelatex main_final4.tex
%  Desde final3 (observaciones del Ing. Miranda, 2026-09-25): texto en Times
%  New Roman de 12 puntos, justificado y con particion silabica.
%  Lo que cambia en final4 (tribunal, defensa de borrador): los rotulos de
%  tablas y figuras van en REDONDA, el numero de pagina va al PIE con el
%  nombre del capitulo y un filete, y las notas al pie quedan fijas al fondo
%  de la hoja. Todo lo demas de la presentacion sigue APA 7, con las
%  excepciones que declara la Introduccion ("Normas de citacion y presentacion").
% =====================================================================

% --- Motor y fuente de TEXTO: la Times New Roman real, con XeLaTeX ---
% Es la del sistema (C:\Windows\Fonts\times*.ttf), no un clon: solo XeLaTeX (o
% LuaLaTeX) puede usarla, por eso final3 y final4 no compilan con pdflatex. Es el mismo
% mecanismo que la version Arial (docs/notas/2026-09-16_tesis-version-arial.md).
%
% Las FORMULAS siguen en Latin Modern, las mismas de final2 (decision del autor,
% como en la version Arial). lmodern va PRIMERO y a proposito: declara las familias
% MATEMATICAS; recien despues fontspec redefine las familias de TEXTO (rm/sf/tt).
\ifdefined\XeTeXversion\else
  \errmessage{main_final4 se compila con XeLaTeX: latexmk -xelatex main_final4.tex}
\fi
\usepackage{lmodern}
\usepackage{fontspec}
% Ligatures=TeX (-- , --- y las comillas de TeX) va SOLO en las familias de
% texto, nunca en la monoespaciada: en un listado de codigo convierte el "--"
% de Python en una raya y las comillas rectas en curvas.
\setmainfont{Times New Roman}[Ligatures=TeX]
\setsansfont{Arial}[Ligatures=TeX]
% La monoespaciada se queda en Latin Modern Mono y no pasa a una mono de Windows.
% No es cosmetica: medido en la version Arial sobre el Anexo C, con Consolas el
% codigo sale con U+2010 HYPHEN en vez del guion ASCII (y con Ligatures=TeX, ademas
% con U+201D en vez de la comilla recta), asi que lo que el tribunal copia del PDF
% ya no es Python valido. Con Latin Modern Mono no hay un solo caracter espurio.
\setmonofont{Latin Modern Mono}
% es-noindentfirst se retiro al adoptar APA: la norma sangra la primera linea de
% TODO parrafo, incluido el primero que sigue a un titulo.
\usepackage[spanish,es-tabla]{babel}  % "Tabla" (no "Cuadro")
% Los nombres propios no se parten: babel cortaba «Edu-FEM», incluso dentro del
% enunciado de la hipotesis (final4), y «SA-FE» en la bibliografia. Las palabras
% inglesas se parten como en ingles, no con los patrones del espanol («softwa-re»,
% «Cambrid-ge»).
\hyphenation{EduFEM SAFE soft-ware Cam-bridge}

% --- Parche: babel-spanish rompe los \caption con formulas bajo XeLaTeX ---
% Sintoma exacto: "! Argument of \es@a has an extra }." al procesar un
% \caption{... $D_{\max}$.}, que abortaba la compilacion con XeLaTeX (visto
% primero en la version Arial; final3 compila con XeLaTeX y lo necesita).
% Causa: babel-spanish redefine \max, \min, \lim e \inf para que impriman "max",
% "min", "lim" e "inf" ACENTUADOS, que es lo correcto en espanol. Para poner ese
% acento elige entre dos implementaciones segun el motor, y la que reserva para
% XeLaTeX (\es@op@ac@TU, activada porque ahi existe \Umathchar) no sobrevive a la
% expansion que hace \caption al escribir el rotulo en el archivo .lof.
% La rama base -la que usa pdflatex, o sea la que ya produce main.pdf- si.
% Verificado: con este \let, el texto extraido del PDF es identico al de
% pdflatex ("max" "min" "lim" "inf" con tilde); sin el, no compila. La
% alternativa \unaccentedoperators tambien evita el error pero deja los
% operadores SIN acento, apartandose de main.pdf.
\ifdefined\XeTeXversion
  \makeatletter
  \@ifundefined{es@op@ac@base}{}{\let\es@op@ac\es@op@ac@base}
  \makeatother
\fi
\usepackage{csquotes}            % requerido por biblatex

% babel-spanish redefine \% para anteponer un espacio fino, y esa version rompe en
% modo matematico ("Incompatible glue units"). Se restablece un \% simple, valido en
% texto y en formulas (los capitulos ya escriben el espacio fino con \, donde hace falta).
\AtBeginDocument{\DeclareRobustCommand{\%}{\char37\relax}}

% Guion bajo que admite corte de linea justo despues de el. Permite que los
% identificadores de codigo largos (p.ej. generate_structured_quad_mesh,
% boundary_conditions.csv) se ajusten dentro de las celdas p{} estrechas de las
% tablas de anexos sin desbordar el margen. Es puramente aditivo: solo agrega
% oportunidades de corte, nunca fuerza un salto.
\AtBeginDocument{\DeclareRobustCommand{\_}{\textunderscore\allowbreak}}

% =====================================================================
%  PRESENTACION DEL DOCUMENTO: formato APA 7.a edicion, con cinco excepciones
% ---------------------------------------------------------------------
%  Las CITAS y la lista de REFERENCIAS no son APA: siguen Vancouver
%  (numericas, por orden de aparicion). Ver el bloque de biblatex mas
%  abajo y tesis/normas/guia_vancouver.pdf. Lo que sigue es solo la
%  presentacion: margenes, fuente, interlineado, sangria, titulos,
%  rotulos de tablas y figuras y numeracion de pagina.
%  Excepciones a APA, declaradas en la Introduccion: capitulos y secciones
%  numerados, papel A4, (desde final3) texto JUSTIFICADO con particion
%  silabica y (desde final4, pedido del tribunal) el titulo de tablas y
%  figuras en REDONDA y el numero de pagina AL PIE con el nombre del
%  capitulo. La fuente, Times New Roman 12, es una de las que APA admite.
%  Que se aplico de APA y por que: docs/notas/2026-09-10_apa-presentacion.md
% =====================================================================

% --- Geometria, interlineado y parrafos (APA 7) ---
% APA pide 2,54 cm (1 pulgada) en los cuatro lados. Se conserva A4 en vez de
% carta: es el papel de la region y el de la caratula ya impresa.
% headheight: heredado de final3, cuando el numero de pagina iba en el encabezado
% (fancyhdr necesitaba 14,5 pt a 12 pt). Hoy el encabezado va vacio y no estorba.
% footskip (final4): el pie vive en el margen inferior, SIN includefoot, asi que la
% caja de texto y la paginacion no cambian. Es la distancia de la ultima linea base
% de la caja a la del pie: con 1,27 cm el numero queda a 1,27 cm del borde de la
% hoja y el filete del pie, a 0,7 cm de la caja (con los 30 pt de report quedaba
% demasiado cerca de una nota al pie fijada al fondo).
\usepackage[a4paper,margin=2.54cm,headheight=14.5pt,footskip=1.27cm]{geometry}
\usepackage{setspace}
% APA: interlineado doble en todo el documento. Medido sobre esta tesis: a doble
% son 143 hojas y a 1,5 son 120 (el original, antes de APA, eran 131). El salto
% es chico porque el margen de 2,54 cm ensancha la caja de texto y compensa.
% Para pasar a 1,5, cambiar esta linea por \onehalfspacing.
\doublespacing
\setlength{\parskip}{0pt}        % APA: sin espacio extra antes ni despues del parrafo
\setlength{\parindent}{1.27cm}   % APA: sangria de primera linea de 1/2 pulgada
% Viudas, huerfanas y palabras partidas entre dos hojas (final4). Con \raggedbottom, el
% \parskip rigido y un interlineado que no estira, una hoja corta no tiene con que
% estirarse y TeX no pondera las penalizaciones por defecto (150 y 100): cortaba donde
% caia y dejaba solo arriba el ultimo renglon de un parrafo («miento de calculo.», p. 5),
% solo abajo el primero o una palabra partida entre dos hojas («de-» | «fecto»,
% pp. 108-109). Con 10000 esos cortes quedan prohibidos y la hoja termina un renglon
% antes. LaTeX copia \clubpenalty en \@clubpenalty al empezar el documento, asi que
% rige tambien despues de cada titulo y de cada \item.
\widowpenalty=10000
\clubpenalty=10000
\brokenpenalty=10000

% --- Notas al pie (final4) ---
% [bottom]: report en una sola cara compone con \raggedbottom, y en una pagina corta
% la nota subia hasta quedar pegada al texto (en final3, p. 65 a 3,2 cm del fondo de
% la caja y p. 67 a 1,5 cm). Con [bottom] la nota va siempre al fondo de la caja.
% [hang]: la llamada al ras del margen y todos los renglones de la nota alineados a
% su derecha. footmisc detecta setspace si se carga despues de el: la nota sigue a
% espacio sencillo. Va antes de hyperref. La URL de las notas sale en la letra del
% texto (solo en las notas: la bibliografia conserva su formato).
\usepackage[bottom,hang]{footmisc}
\setlength{\footnotemargin}{1em}
\renewcommand{\footnotelayout}{\urlstyle{same}}

% Excepciones al doble espacio que la propia APA admite: el cuerpo de las tablas
% y el texto dentro de figuras "puede utilizar interlineado sencillo, 1,5 o doble,
% el que sea mas efectivo". Se agregan los listados de codigo por el mismo motivo:
% a doble espacio un fragmento de Python se vuelve ilegible.
%
% TEXTO JUSTIFICADO (final3, pedido del tribunal). Hasta final2 el texto iba a la
% izquierda, sin justificar y sin guiones, como pide APA; final3 lo justifica y
% por eso DEVUELVE la particion silabica: sin guiones, un renglon justificado solo
% puede estirar los espacios, y deja huecos o se mete en el margen. Los patrones
% de particion son los de babel-spanish. ragged2e se sigue cargando, SIN la
% opcion [document], porque los titulos de seccion usan \RaggedRight y los
% indices \justifying.
\usepackage{ragged2e}
% Red de seguridad del justificado: si un parrafo no encuentra cortes buenos,
% TeX puede estirar un poco mas los espacios antes de dejar un renglon desbordado
% (URL, identificadores de codigo, formulas en linea).
\setlength{\emergencystretch}{1.5em}

\usepackage{etoolbox}
\AtBeginEnvironment{tabular}{\singlespacing}
\AtBeginEnvironment{tabularx}{\singlespacing}
\AtBeginEnvironment{longtable}{\singlespacing}
\AtBeginEnvironment{lstlisting}{\singlespacing}

% --- Las celdas p{} de las tablas van a la izquierda, sin sangria y con guiones ---
% El texto corrido va justificado, pero una celda angosta justificada abre huecos:
% medido en el primer XeLaTeX de final3, las columnas p{} sin \raggedright de la
% Nomenclatura y de los Anexos B y F dieron 25 renglones flojos (badness 10000 con
% los identificadores de codigo). Las celdas se componen como en final2, a la
% izquierda, que es lo que APA permite dentro de una tabla ("como resulte mas
% efectivo"). \@arrayparboxrestore lo usan las celdas p{} y \parbox; en esta tesis
% el unico \parbox es el de \figpend, que se centra por su cuenta.
\makeatletter
\apptocmd{\@arrayparboxrestore}{%
  \RaggedRight\parindent\z@
  \hyphenpenalty=50 \exhyphenpenalty=50 % los valores por defecto de LaTeX
}{}{\PackageError{tesis}{No se pudo parchear \string\@arrayparboxrestore}{}}

% --- Los indices (contenido, figuras, tablas) van justificados y con guiones ---
% Heredado de final2, donde el resto del texto iba sin justificar: sus renglones
% terminan en un relleno de puntos y el numero de pagina, y sin guiones el titulo
% de una entrada larga se metia en el margen (con Arial, 8 pt). En final3 todo el
% texto ya va justificado y con guiones; el parche queda como guarda inocua.
\let\tesis@starttoc\@starttoc
\renewcommand{\@starttoc}[1]{\begingroup\justifying
  \hyphenpenalty=50 \exhyphenpenalty=50
  \tesis@starttoc{#1}\endgroup}
\makeatother

% --- Matematica ---
\usepackage{amsmath,amssymb,amsfonts}
\usepackage{bm}                  % \bm para negrita matematica (matrices)

% Operador de ensamblaje del MEF: la "A grande" estandar (no el wedge logico).
% Con \DeclareMathOperator* el subindice (el indice de elemento) se ubica debajo.
\DeclareMathOperator*{\Ensamblaje}{\text{\Large\textsf{A}}}

% --- Figuras y tablas ---
\usepackage{graphicx}
\graphicspath{{capitulos_final4/figuras/}{figuras/}}   % final4: las figuras rehechas tienen prioridad
\usepackage{booktabs}
\usepackage{array}
\usepackage{longtable}
\usepackage{multirow}
\usepackage{float}
\usepackage{placeins}  % \FloatBarrier: la Figura CR.2 queda antes de «Recomendaciones»
\usepackage{needspace} % \needspace: titulos de grupo de las Recomendaciones sin quedar solos al pie
\usepackage{caption}
% --- Rotulos de tablas y figuras (APA 7, salvo la cursiva) ---
% APA: numero en negrita en su propio renglon y titulo debajo, encima del
% contenido; la nota va debajo del contenido (comando \notaapa). final3: el titulo
% va justificado, como el resto del texto (en final2 iba a la izquierda).
% final4: APA pide el titulo en cursiva; va en REDONDA por pedido del tribunal
% (defensa de borrador). Excepcion declarada en la Introduccion. Un \emph dentro
% del titulo (p. ej. un termino en ingles) sale ahora en cursiva, como en el texto.
\captionsetup{labelfont=bf,        % "Figura 1.1" en negrita
              textfont=up,         % titulo en redonda (final4)
              labelsep=newline,    % titulo en el renglon siguiente
              singlelinecheck=false,
              justification=justified,
              font=normalsize}

% APA coloca numero y titulo ARRIBA de la figura, no debajo. plaintop lo resuelve
% sin tocar los 29 entornos figure de los capitulos, que traen el \caption al final.
\floatstyle{plaintop}
\restylefloat{figure}

% Nota al pie de una tabla o figura, en el formato de APA: "Nota." en cursiva.
% Uso: dentro del float, despues del contenido.  \notaapa{Elaboracion propia.}
\newcommand{\notaapa}[1]{%
  \par\vspace{0.4em}%
  {\justifying\footnotesize\textit{Nota.} #1\par}}
\usepackage{pdfpages}            % embeber PDF externo (Anexo de validacion SAP2000)

% --- Color, enlaces y tipografia fina ---
\usepackage{xcolor}
\definecolor{tesisAzul}{HTML}{0D6EFD}
\usepackage{microtype}
\usepackage[hidelinks]{hyperref}
\hypersetup{
  pdftitle={\tituloTesis},
  pdfauthor={\autorTesis},
  pdfsubject={Tesis de licenciatura en Ingeniería Civil, Universidad Autónoma Tomás Frías},
  pdfkeywords={método de elementos finitos, software educativo, trazabilidad, verificación y validación, memoria de cálculo},
  bookmarksnumbered=true,
}
% \autoref{anexo:...} debe imprimir "Anexo X" (el rotulo de los capitulos del apendice)
% tambien ANTES de \appendix (la Introduccion los enumera). babel-spanish restablece
% los nombres de hyperref en cada seleccion de idioma, por eso se agrega a \extrasspanish.
\addto\extrasspanish{\renewcommand{\appendixautorefname}{Anexo}}

% --- Listados de codigo fuente (Anexo C) ---
% El paquete listings con un mapeo `literate` para los simbolos unicode que
% aparecen en los comentarios del codigo (xi, eta, derivadas parciales, letras
% griegas), de modo que los fragmentos reales del proyecto se reproduzcan fieles.
\usepackage{listings}
\definecolor{lstKeyword}{HTML}{0D6EFD}
\definecolor{lstComment}{HTML}{6C757D}
\definecolor{lstString}{HTML}{198754}
\definecolor{lstBg}{HTML}{F6F8FA}
\lstdefinestyle{pyedufem}{
  language=Python,
  basicstyle=\ttfamily\footnotesize,
  keywordstyle=\color{lstKeyword}\bfseries,
  commentstyle=\color{lstComment}\itshape,
  stringstyle=\color{lstString},
  backgroundcolor=\color{lstBg},
  numbers=left, numberstyle=\tiny\color{lstComment}, numbersep=8pt,
  % Los numeros de linea cuelgan a la izquierda del marco: sin este margen
  % invadian el margen APA de 2,54 cm.
  xleftmargin=1.6em, xrightmargin=0.2em,
  frame=single, framesep=4pt, rulecolor=\color{lstComment},
  breaklines=true, breakatwhitespace=false, showstringspaces=false, upquote=true,
  tabsize=4, columns=fullflexible, keepspaces=true,
  extendedchars=true, inputencoding=utf8,
  literate=%
    {ξ}{{$\xi$}}1 {η}{{$\eta$}}1 {∂}{{$\partial$}}1 {∈}{{$\in$}}1
    {ε}{{$\varepsilon$}}1 {γ}{{$\gamma$}}1 {σ}{{$\sigma$}}1 {τ}{{$\tau$}}1
    {ν}{{$\nu$}}1 {λ}{{$\lambda$}}1 {μ}{{$\mu$}}1 {θ}{{$\theta$}}1 {π}{{$\pi$}}1
    {·}{{$\cdot$}}1 {×}{{$\times$}}1 {→}{{$\rightarrow$}}1 {∇}{{$\nabla$}}1
    {²}{{\textsuperscript{2}}}1 {³}{{\textsuperscript{3}}}1
    {≤}{{$\leq$}}1 {≥}{{$\geq$}}1 {≈}{{$\approx$}}1 {±}{{$\pm$}}1 {√}{{$\surd$}}1
    {—}{{---}}1 {–}{{--}}1 {°}{{$^{\circ}$}}1
    {á}{{\'a}}1 {é}{{\'e}}1 {í}{{\'i}}1 {ó}{{\'o}}1 {ú}{{\'u}}1
    {Á}{{\'A}}1 {É}{{\'E}}1 {Í}{{\'I}}1 {Ó}{{\'O}}1 {Ú}{{\'U}}1
    {ñ}{{\~n}}1 {Ñ}{{\~N}}1 {ü}{{\"u}}1 {¿}{{?`}}1 {¡}{{!`}}1
}
\lstset{style=pyedufem}

% --- XeLaTeX: los simbolos por encima de U+00FF en el codigo (final3) ---
% Bajo XeLaTeX, listings no aplica `literate` a los caracteres por encima de 255:
% los pasa tal cual a la fuente monoespaciada, y Latin Modern Mono no tiene ∂ ni
% las griegas (medido en el primer XeLaTeX de final3: 34 "Missing character" en
% el Anexo C, que habrian salido en blanco). Se mapean aqui al simbolo matematico,
% que es exactamente lo que `literate` hacia con pdflatex. Va DESPUES del estilo
% de listings: si esos caracteres fueran activos al leer `literate`, su lista se
% romperia. Solo aparecen en los listados; en el texto, la tesis escribe las
% griegas en modo matematico.
\usepackage{newunicodechar}
\newunicodechar{∂}{\ensuremath{\partial}}
\newunicodechar{∇}{\ensuremath{\nabla}}
\newunicodechar{∈}{\ensuremath{\in}}
\newunicodechar{ξ}{\ensuremath{\xi}}
\newunicodechar{η}{\ensuremath{\eta}}
\newunicodechar{ε}{\ensuremath{\varepsilon}}
\newunicodechar{γ}{\ensuremath{\gamma}}
\newunicodechar{σ}{\ensuremath{\sigma}}
\newunicodechar{τ}{\ensuremath{\tau}}
\newunicodechar{ν}{\ensuremath{\nu}}
\newunicodechar{λ}{\ensuremath{\lambda}}
\newunicodechar{μ}{\ensuremath{\mu}}
\newunicodechar{θ}{\ensuremath{\theta}}
\newunicodechar{π}{\ensuremath{\pi}}
\newunicodechar{≤}{\ensuremath{\leq}}
\newunicodechar{≥}{\ensuremath{\geq}}
\newunicodechar{≈}{\ensuremath{\approx}}
\newunicodechar{√}{\ensuremath{\surd}}

% --- Pie de pagina (final4): numero al pie, con el capitulo y un filete ---
% APA (trabajo de estudiante) pone UNICAMENTE el numero de pagina, arriba a la
% derecha. final4 lo lleva al PIE por pedido del tribunal (excepcion declarada en
% la Introduccion): a la izquierda, el nombre corto del capitulo vigente; a la
% derecha, el numero; encima, un filete fino. El encabezado queda vacio.
% La hoja de apertura de cada capitulo (estilo plain) lleva solo el numero, en el
% mismo sitio y sin filete: el titulo del capitulo ya esta en ella.
% \leftmark es la marca del capitulo. La pone \chaptermark (redefinido abajo, con
% punto medio), salvo en los capitulos sin numero, que fijan la suya con \markboth,
% y en los tres capitulos cuyo titulo es el objetivo especifico completo: ese
% titulo no cabe en el pie, y cada uno llama a \chaptermark con un nombre corto
% justo despues de su \label. \nouppercase anula las mayusculas que report y
% biblatex ponen en sus marcas (indices, bibliografia).
\usepackage{fancyhdr}
\pagestyle{fancy}
\fancyhf{}
\fancyfoot[L]{\small\nouppercase{\leftmark}}
\fancyfoot[R]{\thepage}
\renewcommand{\headrulewidth}{0pt}
\renewcommand{\footrulewidth}{0.4pt}
\renewcommand{\footruleskip}{6pt}
\fancypagestyle{plain}{% las hojas de apertura de capitulo usan "plain" en report
  \fancyhf{}%
  \fancyfoot[R]{\thepage}%
  \renewcommand{\headrulewidth}{0pt}%
  \renewcommand{\footrulewidth}{0pt}%
}
% La Nomenclatura (una longtable larga) perdia su marca en las hojas de
% continuacion, que imprimian la del indice anterior: estilo con el nombre fijo.
% Se activa SOLO dentro de un grupo (00b_nomenclatura.tex): fuera de el
% reescribiria los campos del estilo fancy para el resto del documento.
\fancypagestyle{nomenclatura}{%
  \fancyhf{}%
  \fancyfoot[L]{\small Nomenclatura}%
  \fancyfoot[R]{\thepage}%
  \renewcommand{\headrulewidth}{0pt}%
  \renewcommand{\footrulewidth}{0.4pt}%
}
% fancyhdr define \chaptermark una sola vez, al cargarse; esta version la reemplaza:
% "Capitulo 2 · Desarrollo de EduFEM", "Anexo A · Manual de instalacion".
\makeatletter
\renewcommand{\chaptermark}[1]{%
  \markboth{\@chapapp~\thechapter\ \textperiodcentered\ #1}{}}
\makeatother

% --- Jerarquia de titulos en formato APA 7 (conservando la numeracion) ---
% APA define cinco niveles, todos en negrita y con cada palabra en mayuscula.
% APA ademas pide no numerar los titulos; aca SI se numeran (decision del autor):
% la tesis tiene ~50 referencias cruzadas del tipo \autoref{cap:...} y el tribunal
% espera capitulos numerados. Lo que se adopta es la tipografia de cada nivel.
%
%   Nivel 1 (\chapter)       centrado, negrita
%   Nivel 2 (\section)       izquierda, negrita
%   Nivel 3 (\subsection)    izquierda, negrita cursiva
%   Nivel 4 (\subsubsection) sangrado, negrita, punto final, texto en la misma linea
%   Nivel 5 (\paragraph)     sangrado, negrita cursiva, punto final, texto seguido
%
% El cuerpo es el mismo del texto (APA: "utiliza el mismo tamano de fuente en
% todo el documento, incluso en los titulos"), por eso van en \normalsize.
\usepackage{titlesec}

\titleformat{\chapter}[display]
  {\normalfont\normalsize\bfseries\centering}
  {\chaptertitlename\ \thechapter}{0.6\baselineskip}{}
\titlespacing*{\chapter}{0pt}{0pt}{\baselineskip}

% [block] y no la forma "hang" por defecto: con hang, un titulo largo -"1.8
% Fenomenos numericos en elementos de bajo orden: bloqueo y modos espurios"- no
% envuelve y se mete 0,9 cm en el margen derecho. En bloque, parte a dos
% renglones. El \RaggedRight explicito evita que el titulo salga justificado.
\titleformat{\section}[block]
  {\normalfont\normalsize\bfseries\RaggedRight}{\thesection}{1em}{}
\titlespacing*{\section}{0pt}{\baselineskip}{0pt}

\titleformat{\subsection}
  {\normalfont\normalsize\bfseries\itshape}{\thesubsection}{1em}{}
\titlespacing*{\subsection}{0pt}{\baselineskip}{0pt}

\titleformat{\subsubsection}[runin]
  {\normalfont\normalsize\bfseries}{\thesubsubsection}{1em}{}[.]
\titlespacing*{\subsubsection}{\parindent}{\baselineskip}{0.5em}

\titleformat{\paragraph}[runin]
  {\normalfont\normalsize\bfseries\itshape}{}{0pt}{}[.]
\titlespacing*{\paragraph}{\parindent}{\baselineskip}{0.5em}

% --- Sangria francesa de la lista de referencias (APA: 1,27 cm) ---
% Aplica a la lista de biblatex, que es Vancouver en su contenido pero se
% presenta con la sangria que pide APA, como el resto del documento.
\AtBeginDocument{\setlength{\bibhang}{1.27cm}}

% --- Bibliografia estilo VANCOUVER (biblatex + biber) ---
% Numeracion por orden de aparicion (sorting=none), citas numericas compactas.
% numeric-comp + sorting=none + el formato de nombres de abajo reproduce Vancouver/NLM.
\usepackage[backend=biber,    % obligatorio: biber, NO bibtex
            style=numeric-comp,  % numerico con rangos comprimidos: [2-5]
            sorting=none,        % numeracion por orden de aparicion (Vancouver)
            sortcites=true,      % ordena numeros dentro de una cita: [3,1] -> [1,3]
            giveninits=true,     % nombres como iniciales
            terseinits=true,     % iniciales pegadas y sin puntos: "Bathe KJ", "Ryu HR".
                                 % Fija bibinitperiod, bibinitdelim Y bibinithyphendelim;
                                 % este ultimo (por defecto ".-") es el que hacia salir
                                 % "Bathe K.-J." en los nombres compuestos.
            uniquename=false,
            uniquelist=false,
            maxbibnames=6,       % hasta 6 autores en la lista...
            minbibnames=6,       % ...y si hay mas, corta con "et al." (Vancouver/NLM)
            doi=true,            % mostrar DOI en articulos
            isbn=false,
            url=false]{biblatex}

% Formato de nombres estilo Vancouver/NLM: "Salari K" (apellido + iniciales,
% sin coma, sin puntos ni espacios entre iniciales) y sin "y"/"and" final.
% (No hace falta redefinir "sortname": con sorting=none nunca se usa.)
\DeclareNameAlias{default}{family-given}
\renewcommand*{\revsdnamepunct}{}        % sin coma entre apellido e iniciales
% (bibinitperiod y bibinitdelim los deja vacios la opcion terseinits de arriba,
%  junto con bibinithyphendelim, que estos dos \renewcommand no alcanzaban.)
% El ultimo autor se separa con "y" (p.ej. "Salari K y Knupp P"): es la convencion
% de babel-spanish y la natural en un documento en espanol. El NLM ingles usa coma;
% para forzarla habria que sobrescribir el delimitador de babel, lo que es fragil.

\addbibresource{capitulos_final4/referencias.bib}   % final4: copia propia de la bibliografia

% Citas multiples con localizador (\autocites[p.~a]{k1}[p.~b]{k2}): el delimitador
% por defecto (coma) era ambiguo, "[20, pp. 21, 26, 19, pp. 219-222]". Se redeclara
% SOLO el comando multiple \parencites (al que \autocites se mapea con autocite=inline)
% con punto y coma, "[20, pp. 21, 26; 19, pp. 219-222]", y se deja \multicitedelim
% en coma para que \autocite{k1,k2} siga imprimiendo "[4, 14]" como pide Vancouver.
\DeclareMultiCiteCommand{\parencites}[\mkbibbrackets]{\parencite}{\addsemicolon\space}
\ExecuteBibliographyOptions{autocite=inline}

% --- Comando para figuras todavia inexistentes (placeholder compilable) ---
% Uso: dentro de un entorno figure, \figpend{descripcion de la figura}.
% Cuando exista la imagen real, reemplazar por \includegraphics{...}.
\newcommand{\figpend}[1]{%
  \begin{center}
  \setlength{\fboxsep}{10pt}%
  \fcolorbox{tesisAzul}{white}{%
    \parbox[c][3.4cm][c]{0.82\linewidth}{\centering
      \textbf{[Figura pendiente]}\\[4pt]\small\itshape #1}}%
  \end{center}}

% Etiqueta de "pendiente" para texto que falta completar.
\newcommand{\pendiente}[1]{\textcolor{red}{\textbf{[PENDIENTE: #1]}}}

% final1: anchos de los indices (H-219, H-339, H-257)
% El numero de pagina se pegaba al texto en las entradas largas y el numero de figura de dos
% digitos («B.10.») al titulo: se ensancha la caja del numero y el margen derecho de las lineas.
\makeatletter
\renewcommand{\@pnumwidth}{2.2em}
\renewcommand{\@tocrmarg}{3.2em}
\renewcommand*\l@figure{\@dottedtocline{1}{1.5em}{3.0em}}
\renewcommand*\l@table{\@dottedtocline{1}{1.5em}{3.0em}}
\makeatother
